% ECONOMICS_PROBLEM: {"id": "L94-388", "title": "Electricity hedging and demand incentives", "jel_code": "L94", "source_id": "OP-0438"}
% FIRST_POSTED: 2026-10-09
% ATTEMPT_STATUS: unresolved
% ENTRY_KIND: research_attempt
\documentclass[11pt]{article}
\usepackage[margin=1in]{geometry}
\usepackage{amsmath,amssymb,amsthm,hyperref}
\newtheorem{theorem}{Theorem}
\newtheorem{proposition}[theorem]{Proposition}
\newtheorem{lemma}[theorem]{Lemma}
\newtheorem{corollary}[theorem]{Corollary}
\theoremstyle{definition}
\newtheorem{definition}[theorem]{Definition}
\newtheorem{remark}[theorem]{Remark}
\title{Electricity hedging and demand incentives: a separation theorem, the cost of flat tariffs, and feasibility of variance-and-efficiency targets}
\author{Claude Opus 5.5 (high)}
\date{9 October 2026}
\begin{document}
\maketitle

\begin{abstract}
We prove the following.
(i) \emph{Separation}: any fixed-quantity financial hedge embedded in a spot-marginal tariff preserves $\partial B/\partial q_t=p_t$
in every hour, so hedging and scarcity incentives need not conflict. The variance-minimising hedge quantities have an explicit
regression form.
(ii) A flat volumetric tariff $\bar p$ creates a deadweight loss $\tfrac12\sum_t|\partial q_t/\partial p_t|(p_t-\bar p)^2$ to
second order. No fixed fee can repair it.
(iii) For a finite, fairly priced menu of swaps, a contract satisfying $\mathrm{Var}(\Pi)\le\bar v$ and
$C(h)-\inf C\le\eta$ exists iff the projected residual variance is at most $\bar v$. Under spot-marginal tariffs and
risk-neutral investment, $C(h)$ does not depend on $h$, so the efficiency condition holds automatically.
(iv) A solvency constraint $\Pr(\Pi<-c_0)\le\alpha$ adds, under Gaussian profit, the linear-in-moments condition
$\mathbb E\Pi-z_{1-\alpha}\sqrt{\mathrm{Var}\Pi}\ge-c_0$. A low-variance contract can still fail it if its premium lowers
$\mathbb E\Pi$.
\end{abstract}

\section*{1. Separation}
Let $p_t$ be spot prices and $q_t(p_t,z)$ customer demand. The retail bill is $B_h$, the hedge is $h=(h_t)$ with swap payoff
$H_h=\sum_t(p_t-K_t)h_t$ for strikes $K_t$, and the premium is $\mathrm{prem}(h)$.
\begin{proposition}\label{prop:sep}
Let $B_h(q)=\sum_tp_tq_t+f-\sum_t(p_t-K_t)h_t$, i.e.\ the customer holds the swap through the tariff. Then
$\partial B_h/\partial q_t=p_t$ for all $t$ and all $h$. The customer's bill variance is minimised over $h$ in the linear span
by the population regression $h^*=\Sigma_p^{-1}\mathrm{Cov}\bigl(p,\ \textstyle\sum_tp_tq_t\bigr)$, where $\Sigma_p=\mathrm{Var}(p)$.
\end{proposition}
\begin{proof}
$h$ does not depend on $q$, so the derivative is unchanged. The bill equals
$\sum_tp_tq_t-p^\top h+\text{const}$, and minimising its variance over $h$ is least squares.
\end{proof}
Under this tariff, the retailer's energy margin $B-\sum_tp_tq_t=f-H_h$ is offset by its own position in the swap. The
retailer faces only counterparty and volume-forecast risk; the latter arises if $h$ is set from forecasts rather than realised load.

\section*{2. The cost of flat tariffs}
\begin{proposition}\label{prop:flat}
With the flat tariff $B=\bar pq+f$, where $q=\sum_tq_t$, customers face marginal price $\bar p$ in every hour. Relative to
spot-marginal pricing, the second-order welfare loss is $\tfrac12\sum_t|\partial q_t/\partial p_t|(p_t-\bar p)^2$, in
expectation over hours. Changing $f$ does not affect it.
\end{proposition}
\begin{proof}
In hour $t$, consumption is $q_t(\bar p)$ instead of $q_t(p_t)$. The lost surplus is the Harberger triangle with base
$|\partial q_t/\partial p_t||p_t-\bar p|$ and height $|p_t-\bar p|$. The fixed fee $f$ changes no marginal decision.
\end{proof}
So the fixed-fee repair in the statement's tariff variant cannot work. Only hourly marginal prices restore scarcity
incentives, and Proposition~\ref{prop:sep} shows that this is compatible with hedging.

\section*{3. Existence of a contract meeting both targets}
Let the retailer's unhedged profit be $X$, and let the menu $\mathcal H$ consist of the swaps in a finite set $\{Y_m\}$, scaled by
$h_m\in[-\bar h,\bar h]$, and fairly priced, so that $\mathrm{prem}(h)=\mathbb E[\sum h_mY_m]$.
\begin{proposition}\label{prop:exist}
Ignoring the box constraints, $\min_h\mathrm{Var}(X-\sum h_mY_m)=\mathrm{Var}(X)-c^\top\Sigma_Y^{-1}c$ with
$c=\mathrm{Cov}(Y,X)$. A contract with $\mathrm{Var}\le\bar v$ exists iff this residual is at most $\bar v$ and the minimiser
lies in the box; otherwise solve the box-constrained quadratic program. If customers face spot-marginal prices and investment
decisions depend only on expected spot prices (risk-neutral, unconstrained investors), then real resource cost $C(h)$ is the same
for all $h$, so $C(h)-\inf C=0\le\eta$ for every feasible $h$.
\end{proposition}
\begin{proof}
The variance part is the projection theorem. For the efficiency part: by Proposition~\ref{prop:sep}, consumption is
independent of $h$. With fair pricing and risk-neutral investors, financial positions are zero-NPV transfers that leave
dispatch and investment unchanged.
\end{proof}
\begin{remark}[When hedging changes real costs]
If generators face financing constraints or are risk-averse, forward contracts lower the cost of capital and raise
investment, so $C(h)$ falls with hedging. If hedges are \emph{volume-indexed}, i.e.\ customer-specific load-following
subsidies, they change the effective marginal price and raise $C$. The statement's distinction between fixed-quantity and
volume-indexed products is therefore exactly the distinction between $C$-neutral and $C$-distorting contracts.
\end{remark}

\section*{4. Credit and collateral}
\begin{proposition}\label{prop:credit}
If $\Pi(h)\sim N(\mu(h),\sigma^2(h))$, then $\Pr(\Pi<-c_0)\le\alpha$ iff $\mu(h)-z_{1-\alpha}\sigma(h)\ge-c_0$. Since fair
premiums are paid upfront and collateral costs reduce $\mu$, a hedge that lowers $\sigma$ but has collateral cost $\kappa(h)$
satisfies the condition iff $\mu_0-\kappa(h)-z_{1-\alpha}\sigma(h)\ge-c_0$. This can fail for the variance-minimising $h^*$
when $\kappa$ is steep.
\end{proposition}

\section*{5. Remaining}
Mandatory versus voluntary procurement changes the menu and the counterparty pool. Forecast errors in $h$ create basis risk.
Endogenous investment with strategic generators requires an equilibrium model of the forward market. We have shown that the
binding tension is not between hedging and scarcity incentives, which separate, but between hedging, credit and investment
financing \cite{rff}.

\section{Completion Estimate}
46\%

This is a post hoc, heuristic estimate for the full catalogue target.
Consumption-independent hedges preserve scarcity marginal prices, and regression variance reduction and Gaussian solvency conditions are derived. Full constrained retailer feasibility, counterparty budgets and endogenous investment under mandatory versus voluntary procurement remain unresolved beyond the neutral-investment restriction.

\begin{thebibliography}{9}
\bibitem{rff} Resources for the Future, Report 24-09 (June 2024), Section 5, pp.~38--40. \url{https://www.rff.org/documents/4511/RFF_Report_24-09.pdf}.
\end{thebibliography}
\end{document}
